% GENERATED FILE. Edit canonical lessons, then run npm run generate:books.
% canonical-source-hash: fnv1a64:0f075b33b7b6102e
% canonical-lessons: JA-C104-tsutsumi, JA-C104-fuutou, JA-C104-okurimono, JA-C104-mari, JA-C104-kane

\chapter{Parcel, Envelope, Gift, Ball, Bell}
\label{ch:ja-parcel-envelope-gift-ball-bell}
\hlchaptermodality{104}

\begin{chapteropening}
\textbf{By the end of this chapter:} \emph{I can say a parcel, an envelope, a gift, a ball and a bell.}

Complete the last lesson of chapter 104: I can say a parcel, an envelope, a gift, a ball and a bell.
\end{chapteropening}

\section[\texorpdfstring{\ja{つつみ} (tsutsumi)}{つつみ (tsutsumi)}]{\ja{つつみ} (tsutsumi) — a parcel}

\label{lesson:JA-C104-tsutsumi}



\begin{quote}
Before the new one: say the Japanese for a tray, then the Japanese for a can.
\end{quote}

\subsection*{You'll want to know: \ja{つつみ}}
\textbf{\ja{つつみ}} — \emph{tsutsumi} — \textquotedblleft{}a parcel\textquotedblright{}.

\begin{grammarlens}[title={What you've built}]
One more word for this chapter.
\end{grammarlens}

\subsection*{Your turn}
\begin{itemize}
\raggedright
  \item \emph{Say it:} \emph{tsutsumi}
  \item \emph{Say it:} \emph{tsutsumi}, once more
  \item \emph{Say it:} say \emph{tsutsumi}
  \item \emph{Recall:} say the Japanese for a tray, then the Japanese for a can, then say \emph{tsutsumi} again
\end{itemize}

\begin{tcolorbox}[breakable,colback=teal!4,colframe=teal!35!black,title={Before you move on}]
What does \ja{つつみ} mean? (\textquotedblleft{}A parcel\textquotedblright{}.) Say it once more.
\end{tcolorbox}
\section[\texorpdfstring{\ja{ふうとう} (fuutou)}{ふうとう (fuutou)}]{\ja{ふうとう} (fuutou) — an envelope}

\label{lesson:JA-C104-fuutou}



\begin{quote}
Before the new one: say the Japanese for a can, then the Japanese for a parcel.
\end{quote}

\subsection*{You'll want to know: \ja{ふうとう}}
\textbf{\ja{ふうとう}} — \emph{fuutou} — \textquotedblleft{}an envelope\textquotedblright{}.

\begin{grammarlens}[title={What you've built}]
One more word for this chapter.
\end{grammarlens}

\subsection*{Your turn}
\begin{itemize}
\raggedright
  \item \emph{Say it:} \emph{fuutou}
  \item \emph{Say it:} \emph{fuutou}, once more
  \item \emph{Say it:} say \emph{fuutou}
  \item \emph{Recall:} say the Japanese for a can, then the Japanese for a parcel, then say \emph{fuutou} again
\end{itemize}

\begin{tcolorbox}[breakable,colback=teal!4,colframe=teal!35!black,title={Before you move on}]
What does \ja{ふうとう} mean? (\textquotedblleft{}An envelope\textquotedblright{}.) Say it once more.
\end{tcolorbox}
\section[\texorpdfstring{\ja{おくりもの} (okurimono)}{おくりもの (okurimono)}]{\ja{おくりもの} (okurimono) — a gift}

\label{lesson:JA-C104-okurimono}



\begin{quote}
Before the new one: say the Japanese for a parcel, then the Japanese for an envelope.
\end{quote}

\subsection*{You'll want to know: \ja{おくりもの}}
\textbf{\ja{おくりもの}} — \emph{okurimono} — \textquotedblleft{}a gift\textquotedblright{}.

\begin{grammarlens}[title={What you've built}]
One more word for this chapter.
\end{grammarlens}

\subsection*{Your turn}
\begin{itemize}
\raggedright
  \item \emph{Say it:} \emph{okurimono}
  \item \emph{Say it:} \emph{okurimono}, once more
  \item \emph{Say it:} say \emph{okurimono}
  \item \emph{Recall:} say the Japanese for a parcel, then the Japanese for an envelope, then say \emph{okurimono} again
\end{itemize}

\begin{tcolorbox}[breakable,colback=teal!4,colframe=teal!35!black,title={Before you move on}]
What does \ja{おくりもの} mean? (\textquotedblleft{}A gift\textquotedblright{}.) Say it once more.
\end{tcolorbox}
\section[\texorpdfstring{\ja{まり} (mari)}{まり (mari)}]{\ja{まり} (mari) — a ball}

\label{lesson:JA-C104-mari}



\begin{quote}
Before the new one: say the Japanese for an envelope, then the Japanese for a gift.
\end{quote}

\subsection*{You'll want to know: \ja{まり}}
\textbf{\ja{まり}} — \emph{mari} — \textquotedblleft{}a ball\textquotedblright{}.

\begin{grammarlens}[title={What you've built}]
One more word for this chapter.
\end{grammarlens}

\subsection*{Your turn}
\begin{itemize}
\raggedright
  \item \emph{Say it:} \emph{mari}
  \item \emph{Say it:} \emph{mari}, once more
  \item \emph{Say it:} say \emph{mari}
  \item \emph{Recall:} say the Japanese for an envelope, then the Japanese for a gift, then say \emph{mari} again
\end{itemize}

\begin{tcolorbox}[breakable,colback=teal!4,colframe=teal!35!black,title={Before you move on}]
What does \ja{まり} mean? (\textquotedblleft{}A ball\textquotedblright{}.) Say it once more.
\end{tcolorbox}
\section[\texorpdfstring{\ja{かね} (kane)}{かね (kane)}]{\ja{かね} (kane) — a bell}

\label{lesson:JA-C104-kane}



\begin{quote}
Before the new one: say the Japanese for a gift, then the Japanese for a ball.
\end{quote}

\subsection*{You'll want to know: \ja{かね}}
\textbf{\ja{かね}} — \emph{kane} — \textquotedblleft{}a bell\textquotedblright{}.

\begin{grammarlens}[title={What you've built}]
One more word for this chapter.
\end{grammarlens}

\subsection*{Your turn}
\begin{itemize}
\raggedright
  \item \emph{Say it:} \emph{kane}
  \item \emph{Say it:} \emph{kane}, once more
  \item \emph{Say it:} say \emph{kane}
  \item \emph{Recall:} say the Japanese for a gift, then the Japanese for a ball, then say \emph{kane} again
\end{itemize}

\begin{tcolorbox}[breakable,colback=teal!4,colframe=teal!35!black,title={Before you move on}]
What does \ja{かね} mean? (\textquotedblleft{}A bell\textquotedblright{}.) Say it once more.
\end{tcolorbox}
